\section{Blowups and surface centers}\label{sec:blowup-centers}

To use the cited decomposition theorems, we must check that they preserve all
the data seen by the invariant.  Four issues arise: coefficient fields,
restriction to even cohomology, regularity in \(z\), and faithful embedding of
the center coefficients.
The last issue is the delicate one: Iritani's center Novikov map may have a
kernel, so the independent target coordinates must recover the missing
numerical curve-class data.
\paragraph{The comparison and its lattice.}
Use Iritani's reduced graded source, its ambient and separate center images,
and the independent external coordinates of Section~5.8.2
\cite[Remarks~2.3 and~5.6, Theorem~5.18]{IritaniBlowup}.
Their generic fields embed in a common algebraically closed overfield after
the prescribed exceptional Novikov Laurent extension and ramification.
This does not identify two differently directed Novikov completions.
The comparison is parity preserving; on the even bulk base it and its inverse
restrict to even cohomology. The same maps act on the full supermodule when
odd dimensions are retained below.

Theorem~5.18 gives mutually inverse maps of the specified graded completed
$\C[z]$-modules, intertwining the actual $z$-connections. Their positive
coefficientwise realization is regular in $z$ in both directions. Hence they
preserve the original lattices, and their leading terms conjugate $N$ and
$\im N$. They therefore preserve the elementary modification and conjugate
its residue. A common grading normalization can add a scalar residue;
it leaves the discriminant unchanged, including in the resonant case.

\paragraph{Faithful center coordinates.}
Iritani's center map need not be injective
\cite[(5.15)]{IritaniBlowup}.  After numerical reduction it has the form
\begin{equation}
 Q_Z^d\longmapsto
 Q^{i_*d}q_{\mathrm{exc}}^{-\rho_Z\cdot d/(c-1)}.
 \label{eq:center-novikov-map-atomic}
\end{equation}
Here \(\rho_Z=c_1(N_{Z/Y})\).
The raw map is not the coefficient map used below.  Remark~2.3 gives the
intrinsic divisor-equation-reduced coefficient ring, whose Novikov generators
are the combined symbols
\[
 X_d=Q_Z^d\exp\!\bigl(\sigma^{(2)}\cdot d\bigr);
\]
the divisor coordinates \(\sigma^{(2)}\) are not separate elements of this
ring.  Remark~5.6 places its pushforward under the raw Novikov map over the
image of that reduced ring \cite[Remarks~2.3 and~5.6]{IritaniBlowup}.  In the
external direct sum of Section~5.8.2, by contrast, the variables
\(s_j\in H^*(Z)\) are independent coordinates on the common comparison target,
and the center parameter is identified as
\(\varsigma_j=\varsigma_j^\circ+s_j\)
\cite[(5.47)--(5.48)]{IritaniBlowup}.  The divisor components of these
independent target coordinates retain the curve-class information lost by
\eqref{eq:center-novikov-map-atomic}.

The distinction is modeled by \(x,y\mapsto q\), which identifies two
generators, and \(x\mapsto qe^s,\ y\mapsto qe^t\), which retains independent
characters.  Source coefficients must not already depend on the target
variables \(s,t\).  The comparison uses the following conventions:
\begin{center}
\small
\begin{tabular}{@{}p{.22\linewidth}p{.72\linewidth}@{}}
\toprule
Object & Required data \\
\midrule
Source & Reduced graded coefficients \(\C[[X_d,\sigma']]_{\rm gr}[\sigma^0]\);
no independent source divisor variables. \\
Target & Independent ambient and center-copy coordinates \(t,s_0,\ldots,s_{c-2}\);
all center divisor directions retained. \\
Scalars & The displayed exceptional or fibre Novikov Laurent extensions and
ramification; their \(z\)-free fraction fields in a common algebraic closure. \\
Connections & Graded completed \(z\)-enhancements embedded coefficientwise
in the common field's \(\overline K[[z]]\); comparison and inverse regular. \\
Derivations & Bulk partials, numerical Novikov logarithmic derivatives and
the actual \(z\)-connection, with coordinate-change chain rules. \\
Grading & Only common scalar residue shifts; no independent shifts of residue
eigenlines. \\
\bottomrule
\end{tabular}
\end{center}

\begin{lemma}[Faithful center base change]
\label{lem:faithful-center-base-change}
\coverage{fragment}%
\lean{faithfulCenterBaseChange_targetOnly_injective}%
Let \(R_{Z,\mathrm{src}}^{\mathrm{red}}
=\C[[X_d,\sigma']]_{\mathrm{gr}}[\sigma^0]\) be the coefficient domain
whose graded completed \(z\)-enhancement is Iritani's reduced source ring
\[\C[z][[Q_Ze^{\sigma^{(2)}},\sigma']][\sigma^0].\]
Let \(A_j\) be the
numerical coefficient domain of the external direct sum in Section~5.8.2,
with the independent target coordinates \(t,s_0,\ldots,s_{c-2}\) and with
\(z\) omitted.  The pullback determined by the center Novikov map and
\(\varsigma_j=\varsigma_j^\circ+s_j\) gives a continuous homomorphism
\[
 \iota_j:R_{Z,\mathrm{src}}^{\mathrm{red}}\longrightarrow A_j.
\]
On a combined divisor-equation generator it is
\begin{equation}
 X_d=Q_Z^d\exp\!\bigl(\sigma^{(2)}\cdot d\bigr)
 \longmapsto
 Q^{i_*d}q_{\mathrm{exc}}^{-\rho_Z\cdot d/(c-1)}
 \exp\!\bigl((\varsigma_j^{\circ,(2)}+s_j^{(2)})\cdot d\bigr).
 \label{eq:faithful-center-map}
\end{equation}
This map is injective.  In particular, it induces a field embedding
\[
 \operatorname{Frac}(R_{Z,\mathrm{src}}^{\mathrm{red}})
 \lhook\joinrel\longrightarrow
 \operatorname{Frac}(A_j).
\]
After this field extension, each center summand in Iritani's comparison is the
faithful scalar extension, followed by the formal coordinate change
\(\varsigma_j\), of the intrinsic generic numerical even QDM of \(Z\).
\end{lemma}

\begin{proof}
All completions here use the graded convention of
\cite[Section~2.2]{IritaniBlowup}.  The pullback is defined on the reduced
ring by the String and Divisor Equations, as explained after (5.36) and
(5.47) there.  It need not be defined on the unrestricted bulk power-series
ring.  We prove injectivity using the first nonzero Novikov degree, retaining
the full exponentials in the target divisor coordinates.
Choose integral divisors \(D_1,\ldots,D_\rho\) whose pairings separate
\(N_1(Z)\), and use their dual coordinates in \(s_j^{(2)}\).  The variables
\(s_j\) occur only in the target \(A_j\), not as coefficients in
\(R_{Z,\mathrm{src}}^{\mathrm{red}}\).  This asymmetry is essential.

Filter by an integral ample degree pulled back from \(Y\).  Every fibre of
the untagged monomial map in a bounded piece is finite: the restriction
of an ample divisor from \(Y\) is
ample on \(Z\), and its bounded slice of \(\overline{NE}(Z)\) contains only
finitely many lattice points.  This filtration is respected by the target
map because \(H\cdot i_*d=(i^*H)\cdot d\), and the fixed shifts have
nonnegative Novikov degree.

For a fixed curve class, the coefficient of an element of the reduced ring
is polynomial in the nondivisor and unit coordinates.  Indeed, an element
has finitely many homogeneous degrees and bounded unit degree; every other
even nondivisor coordinate has strictly negative degree.  Thus only
finitely many of its monomials occur at that curve class.  This is where
the reduced, graded source is essential.

Suppose a nonzero source element maps to zero and take its first nonzero
Novikov-degree part.  At that degree only the Novikov-degree-zero part of
\(\varsigma_j^\circ\) enters.  Translation by its nondivisor and unit
components is an injective substitution on the polynomial coefficients,
over the target Laurent coefficient field.  Its divisor component
multiplies each exponential by a nonzero scalar.  Group the terms by their
common untagged target monomial.  On one nonzero group the
remaining relation has the form
\[
 \sum_{d\in F}a_d
 \exp\!\left(\sum_{k=1}^\rho(D_k\cdot d)s_{j,k}\right)=0,
 \qquad a_d\ne0,
\]
over a characteristic-zero coefficient domain, including the remaining bulk
coordinates, that does not contain the variables \(s_{j,k}\).
The exponent vectors are distinct.  Choose an integral
one-parameter restriction on which they remain distinct.  Derivatives of
orders \(0,\ldots,|F|-1\) at the origin give a Vandermonde matrix, hence every
\(a_d\) is zero, a contradiction.  This proves injectivity on the full
reduced source.  We have not truncated the divisor exponentials by bulk
order; doing so could lose the distinction between curve classes.

Remark~2.3 of \cite{IritaniBlowup} gives the intrinsic reduced source, and
Remark~5.6 gives its image under the raw Novikov map~(5.15).
Although that map can be noninjective on the unrestricted source, it retains
the independent \(\sigma^{(2)}\) coordinates in its target.  Its restriction
to the reduced source is injective by the same argument, with zero shift.
Equations~(5.47)--(5.48) then pull that image into the external direct sum over
\(A_j\), where the variables \(s_j\) are independent, and identify the direct
sum with the blowup module.  Associativity of scalar extension therefore
identifies the center summand with scalar extension directly along the
injective composite \(\iota_j\).  Passing to fraction fields is faithfully
flat.  Finally, the inverse of the combined formal coordinate change
\((\tau,\varsigma_0,\ldots,\varsigma_{c-2})\), supplied by
Theorem~5.18(7) and Section~5.8.2, transports the same statement to the
original blowup coordinates.
\end{proof}

\begin{proposition}[Blowup formula]
\label{prop:numerical-blowup}\coverage{absent}%
For $I=I_{\mathrm{exp}}$ or $I_{\mathrm{lat}}$, a smooth center
$Z\subset Y$ of codimension $c\ge2$ satisfies
\begin{equation}
 I(\Bl_ZY)=I(Y)+(c-1)I(Z).
 \label{eq:numerical-blowup}
\end{equation}
For a disconnected center the correction is summed componentwise.
\end{proposition}
\begin{proof}
Iritani's comparison has one ambient and $c-1$ center copies.
Its joint coordinate map has an inverse, so their unit coordinates are
independent. The resultant of two scalar-translated characteristic
polynomials is a nonzero polynomial in the difference of the unit shifts.
Thus their spectra are disjoint generically. Regular lattice transport
preserves each selection condition, and
Lemma~\ref{lem:faithful-center-base-change} identifies each center copy with
its intrinsic generic connection. Contributions therefore add.
Each operation uses its own common field; no successive Laurent expansions
are composed.
\uses{lem:faithful-center-base-change,prop:rank2-rigidity}%
\imports{IritaniBlowup:thm-5.18}%
\end{proof}

\begin{lemma}[The ruled product in full even bulk]
\label{lem:ruled-product}\coverage{absent}%
For a smooth projective curve $C$, the two counts vanish on $C\times\PP^1$.
Its generic even primary factors have rank at most two.
If $g(C)=1$, every rank-two centered operator is zero; if $g(C)>1$,
every rank-two canonical residue discriminant is zero.
\end{lemma}
\begin{proof}
For $g=g(C)\ge1$, use the flat even basis $1,h,p,hp$, with $h$ and $p$
the point classes of the factors and $\int hp=1$. Write the bulk as
$v\,1+ah+bp+chp$, set $Q=q_f e^b$ and $\kappa=2-2g$.
The horizontal Novikov variable remains an independent spectator.
Every genus-zero map to $C$ is constant. The positive-degree moduli space
is $C$ times the corresponding $\PP^1$ moduli space, with no $H^1$
obstruction from the constant factor. A nonzero invariant has exactly one
$h$ factor. An insertion $h$ alone gives a unit on the second factor and
vanishes by the string equation. Otherwise there is one $hp$ insertion;
all other nonunit insertions are $p$, and the dimension equation forces
fiber degree one. Its three-point invariant is one; the divisor equation
handles the remaining $p$ insertions. The full even potential, modulo
terms of degree at most two, is therefore
\[
 F=\tfrac12v^2c+vab+cQ.
\]
Its third derivatives give $h\star h=0$, $h\star p=hp$ and
$p\star p=Q(1+ch)$. The centered Euler element is $\kappa h+2p-chp$.
Put $p'=(1-ch/2)\star p$. Then $(p')^{\star2}=Q$ and
$U=\kappa h+2p'$. After adjoining $\sqrt Q$, the two projectors
$(1\pm p'/\sqrt Q)/2$ have even rank two, with centered operator
$\kappa h$ throughout the full even bulk.
For $g=1$ this is identically zero. For $g>1$ it has rank one; at the
small point the product connection is the tensor connection, whose curve
factor has modified residue $\kappa E_{21}-\tfrac12I$.
Regular rank-one separation of the $\PP^1$ factor changes this residue
only by a scalar. Proposition~\ref{prop:rank2-rigidity} keeps its
discriminant zero in all bulk directions.
For $C=\PP^1$, the small product has four distinct eigenvalues over the
independent Novikov field; its rank-one factors persist in all even bulk
directions by Proposition~\ref{prop:generic-spectral-connection-splitting}.
\uses{prop:rank2-rigidity,lem:cyclic-primary-persistence}%
\imports{Behrend:product}%
\end{proof}

\subsection{Vanishing on centers through dimension two}

\begin{proposition}[Vanishing on low-dimensional centers]
\label{prop:atomic-lowdim}
\coverage{fragment}%
\lean{rankTwoResidueMarker_lowDimensionalOccurrenceNullity}%
\uses{prop:rank2-rigidity}%
\imports{BeauvilleSurfaces:minimal-surface-classification}%
For every smooth projective \(Z\) of dimension at most two,
\(I_{\mathrm{exp}}(Z)=I_{\mathrm{lat}}(Z)=0\).  The same vanishing holds
for every center copy in a quantum comparison.
\end{proposition}

\begin{proof}
It suffices to prove vanishing for \(I_{\mathrm{lat}}\), since
\(I_{\mathrm{exp}}\leq I_{\mathrm{lat}}\).
\emph{Points and curves.}
A point has one rank-one block.  For \(\PP^1\), the relation
\(p\star p=Qe^{t_p}\) gives two rank-one generic blocks.

Let \(C\) have genus \(g\ge1\).  There is no nonconstant genus-zero map to
\(C\), so its even big quantum product is classical.  With
\(\int_Cp=1\) and \(a=2-2g\), its centered even connection is
\[
 z\partial_zY=
 \left[z^{-1}aE_{21}+\diag(1/2,-1/2)\right]Y.
\]
For \(g=1\), \(N=0\).  For \(g>1\), the modification
\(S=\diag(z,1)\) gives residue \(aE_{21}-\tfrac12I\), of discriminant zero.
Its two eigenvalues, and therefore its two exponent classes, are equal.  Thus
every curve has zero invariant.

\smallskip
\noindent\emph{Surfaces with nef canonical class.}
Let \(S\) be a surface with \(K_S\) nef.  For a quantum-product term with
input degree \(b\), non-scalar Euler degree \(a\ge2\), curve class \(d\), and
nonunit even bulk degrees \(e_\nu\ge2\), the dimension axiom gives output
degree
\[
 a+b-2c_1(S)\cdot d+\sum_\nu(e_\nu-2)\ge b+2.
\]
The centered Euler operator strictly raises ordinary degree.  Its only block
is the whole even cohomology, of rank \(b_0+b_2+b_4\ge3\), so its invariant is
zero.

\smallskip
\noindent\emph{Other surfaces.}
A minimal surface with non-nef canonical class is \(\PP^2\) or a ruled
surface \(\PP_C(V)\) \cite[Chapter~VI]{BeauvilleSurfaces}.  The relation
\(H^{\star3}=Q\) makes \(\PP^2\) generically semisimple.  A ruled surface is birational to $C\times\PP^1$: its defining vector
bundle is trivial over the generic point of $C$.
Lemma~\ref{lem:ruled-product} gives zero on that product. Projective surface
factorization has only point centers, whose contributions in
\eqref{eq:numerical-blowup} are zero. Thus every ruled surface has zero count.
Every nonminimal surface is obtained from a minimal one by point blowups;
formula~\eqref{eq:numerical-blowup} adds only zero point terms.  This proves
intrinsic vanishing for every surface.
Lemma~\ref{lem:faithful-center-base-change} identifies every separate center
summand with a faithful scalar extension of the corresponding intrinsic generic
block, and the invariant is unchanged by that extension.
\uses{prop:numerical-blowup,lem:ruled-product,prop:rank2-rigidity,lem:faithful-center-base-change}%
\imports{AKMW:weak-factorization,BeauvilleSurfaces:minimal-surface-classification}%
\end{proof}

